% Local full proofs; source omitted/cited boundaries remain in the note.

\begin{annexitem}{annex:p04-fo-naming}{thm:p04-fo-naming}{Polygon, box and source-clause translations}
\begin{proof}
Finite polygon coding records its length and rational vertex codes.
Rational comparisons decide validity. Evaluation at a rational point,
slopes, extrema and integrations of a polygon use finite rational
arithmetic. For the minimum absolute value on a rational compact
interval, split at its endpoints and polygon vertices; on each
linear piece return zero if it crosses zero, otherwise the smaller
endpoint absolute value. Finite unions are handled componentwise.

To verify density, given continuous \(f\) and \(\epsilon>0\),
uniform continuity gives a rational mesh so that the variation on
each cell is less than \(\epsilon/2\). Choose rational vertex values
within \(\epsilon/2\) of the true values. At any point in a cell,
linear interpolation of those two values differs from \(f\)
by less than \(\epsilon\), by the triangle inequality and convex
weights. Thus names exist. If one stream names both \(f,g\), then
\(\lVert f-g\rVert_\infty\leq2^{1-n}\) for every \(n\), so \(f=g\).

For the native word/code bridges, parse the \(\#\)-separated
rational words using Appendix C1. Decoding their binary integers
and reduced signed fractions gives finite rational values.
Conversely, reduce a finite fraction by Euclid, normalize the
sign and positive denominator, emit an integer word if the
denominator is one, and otherwise emit the source reduced
fraction word. Zero is the single integer word \(0\);
negative zero and unreduced or denominator-one fraction words
are not emitted. These are terminating operations on finite
data. They translate the native finite-word \(\nu_{\mathbb Q}\)
to and from the host numerical rational codes without
identifying their input types.

For native scalar \(\rho_C\) to host scalar
\(\delta_{\mathrm C}\), parse each \(u_n\#\) block and output
its numerical rational code. Taking \(j\to\infty\) in
\(|\bar u_j-\bar u_n|<2^{-n}\) gives
\(|x-\bar u_n|\leq2^{-n}\), the host clause.
Conversely read the host approximation at position \(n+2\),
emit its normalized rational word followed by \(\#\), and
repeat. For \(j>n\) the decoded values have difference
at most \(2^{-j-2}+2^{-n-2}<2^{-n}\) and converge to the
same scalar. This proves both scalar precision/syntax
translation directions.

For native function \(\delta_C\) to
\(\gamma_{\mathrm{pair}}\), parse each complete
\(w_n\text{\textcent}\) block, parse the finite vertex words
inside \(w_n\), and code the same rational polygon numerically.
In the reverse direction decode each valid host polygon code,
normalize its rational vertex words, join them with \(\#\),
and append the cent delimiter. The polygons and their metric
distances are unchanged. Both the strict pairwise clause and
the uniform limit-error clause are preserved.
If the latter is written instead as uniform convergence to
\(f\), taking a metric limit in the pairwise inequality gives
\(d(f,\alpha(w_n))\leq2^{-n}\); conversely those error bounds
already imply that uniform convergence.

A source pairwise name thus has the host error clause.
Conversely, from host \(P_n\) use \(Q_n=P_{n+2}\):
for \(j>n\),
\(\lVert Q_j-Q_n\rVert_\infty
\leq2^{-j-2}+2^{-n-2}<2^{-n}\).
The limit is the same \(f\). Finite rational fraction syntax
converts effectively by the previously supplied rational coding.

From a host name enumerate all pairs of a candidate box and
an approximation index \(k\). Emit a box \((a,b,c,d)\) when
\[
 c<\min_{[a,b]}P_k-2^{-k},
 \qquad \max_{[a,b]}P_k+2^{-k}<d.
\]
Otherwise emit padding. The test is rational and sound.
If \(f([a,b])\subseteq(c,d)\), compactness gives a positive
margin between its extreme values and \(c,d\).
For sufficiently large \(k\) the test passes.
Thus all and only the correct boxes are eventually listed.

For the reverse translation fix output position \(n\) and
\(\epsilon=2^{-n-3}\). Retain enumerated boxes of range width
less than \(\epsilon\). Search for a finite collection whose
closed domain intervals cover \([0,1]\); rational sorting and
gap checks decide this finite covering condition.
Such a collection exists. Around each point continuity provides
a small relative neighbourhood on which \(f\)'s range fits
strictly inside a rational interval of width less than
\(\epsilon\). Choose a rational closed interval whose relative
interior contains the point and lies in that neighbourhood.
Compactness supplies a finite subcover of the interiors.
Every corresponding true box occurs in the valid input name,
so the finite-cover search eventually succeeds.

Sort all chosen interval endpoints with zero and one into
a grid. Every whole grid cell is contained in at least one
chosen domain interval: its interior contains no interval
endpoint, so whichever interval covers an interior point
contains the entire cell, including its endpoints.
At each grid node choose any covering box and assign the
midpoint of that box's range as the polygon value.
This differs from the true node value by less than
\(\epsilon/2\). On one grid cell the true function values
and both true endpoint values lie in a common range interval
of width less than \(\epsilon\).
Convex interpolation and the node errors give
\[
 |P_n(x)-f(x)|<\epsilon+\epsilon/2<2^{-n}.
\]
This constructs the promised host polygon. The search and
finite arithmetic are effective; no modulus has been received
as unmentioned input.

A list of all true boxes exists set-theoretically for each
continuous function, since the rational quadruples are countable.
It determines the function uniquely: different values at a
rational point can be separated by a rational open range and
the singleton compact box at that point; continuity and rational
density then determine the values at every point. This applies
to both the local and global box interpretations.

A global box-name restricts to local boxes by retaining just
ordered domain intervals inside \([0,1]\). All local properties
then occur, so the preceding translation computes the restriction.
Conversely extend \(f\) by clamping the argument to \([0,1]\).
The same extension of \(P_k\) is uniformly within \(2^{-k}\)
on all of \(\mathbb R\). Its extrema on any rational compact
interval are effectively obtained from the clamped interval.
The first box-enumeration test therefore lists a global
box-name of the extension. The omitted vacuous reversed-domain
boxes of the original syntax are decidably generated separately.
These constructions give both directions for the source's
restricted-global representation, as claimed.

More explicitly, native \(\delta_{\mathbb R}\) to host
\(\gamma_{\mathbb R}\) parses each complete
\(u_i\#v_i\#x_i\#y_i\text{\textcent}\) record.
If its decoded compact endpoints are ordered and its range
endpoints strictly ordered, emit the host quadruple index
plus one; otherwise emit padding zero. Every required proper
box is retained. On a valid native name the other true
records have reversed compact endpoints and are vacuous:
a nonempty interval cannot have its nonempty image inside
an empty range.

In the reverse direction read each host entry, skip its
padding, and recode every proper quadruple into normalized
native rational words and separators. Interleave these
records with an effective enumeration of all rational
quadruples having reversed compact endpoints, including
empty-range quadruples. Each finite stage reads a finite
host prefix and emits complete records. There are
infinitely many vacuous records, so the native output is
productive; every proper property from the host name and
every vacuous property is eventually issued.
Thus the output lists exactly the full native
\(\delta_{\mathbb R}\) property set, without treating the
host zero-padding entry as a source record.
This proves both global syntax translations, not just a
one-directional reinterpretation.
\end{proof}
\end{annexitem}

\begin{annexitem}{annex:p04-fo-evaluation}{thm:p04-fo-evaluation}{Evaluation, a computed modulus and maximum}
\begin{proof}
For precision \(n\), read \(P_{n+2}\) and its rational maximum
absolute slope \(L\). Given a scalar Cauchy name \(q_m\) of
\(x\in[0,1]\), choose \(m\) with
\(L2^{-m}\leq2^{-n-2}\).
This rational search terminates, including \(L=0\).
Let \(c=\min(1,\max(0,q_m))\). Clamping does not increase
distance from \(x\), so \(|c-x|\leq2^{-m}\).
The rational output \(P_{n+2}(c)\) has error
\[
 |P_{n+2}(c)-f(x)|
 \leq L2^{-m}+2^{-n-2}\leq2^{-n-1}.
\]
Every output is therefore a valid scalar approximation.

For a modulus use the same \(P_{n+2}\) and choose \(\mu(n)\)
with \(L2^{-\mu(n)}\leq2^{-n-1}\).
Then for \(|x-y|\leq2^{-\mu(n)}\),
\[
 |f(x)-f(y)|
 \leq2^{-n-2}+L|x-y|+2^{-n-2}\leq2^{-n}.
\]
Each integer \(\mu(n)\) is effectively obtained from the
finite input polygon. Finally extrema of a polygon are its
vertex extrema, and
\[
 |\max P_n-\max f|\leq\lVert P_n-f\rVert_\infty\leq2^{-n}.
\]
The rational maxima thus give a scalar Cauchy name.
\end{proof}
\end{annexitem}

\begin{annexitem}{annex:p04-fo-integration}{thm:p04-fo-integration}{Area error, primitive output and arbitrary limits}
\begin{proof}
The integral of a rational polygon is the finite trapezoid sum
\(\sum_j(t_{j+1}-t_j)(y_j+y_{j+1})/2\), a computable rational.
Since the interval has length one,
\[
 \left|\int_0^1P_n-\int_0^1f\right|
 \leq\int_0^1|P_n-f|\leq2^{-n}.
\]
This proves the scalar algorithm.

For function output let \(M=\lVert P_0\rVert_\infty+1\).
Then \(\lVert f\rVert_\infty\leq M\) and
\(\lVert P_k\rVert_\infty\leq M+1=B\).
At output \(n\), take \(k=n+2\) and a dyadic mesh
\(h=2^{-r}\) with \(Bh\leq2^{-n-2}\).
The values of \(JP_k\) at mesh points are exactly rational:
split at the finitely many polygon vertices and integrate
each linear piece. Interpolate these values linearly.
The function \(JP_k\) is \(B\)-Lipschitz, so on one cell its
distance from the linear interpolant is at most \(Bh\),
using convex endpoint weights. Also
\(\lVert Jf-JP_k\rVert_\infty\leq2^{-k}\).
The output error is at most \(2^{-n-1}\).
Both searches and the output polygon code are effective.

For \(a,b\in[0,1]\), evaluate the computed \(Jf\) at each
and subtract: \(Jf(b)-Jf(a)=\int_a^bf\).
This equation includes equal and reversed limits, without
an equality test.

For global \(f\), use the zeroth scalar approximations of
the input limits to choose a positive integer \(M\) strictly
larger than both \(|a|\) and \(|b|\); their error at position
zero is at most one, so this is effective finite rational
data. Filter the global boxes inside \([-M,M]\) and
rescale their rational endpoints to obtain a local box-name
of \(g(t)=f(-M+2Mt)\). The naming theorem gives its polygon
name. Scalar arithmetic computes
\(\alpha=(a+M)/(2M)\), \(\beta=(b+M)/(2M)\) in \([0,1]\).
The already proved algorithms compute
\[
 2M\bigl(Jg(\beta)-Jg(\alpha)\bigr)=\int_a^bf.
\]
All name-dependent choices of \(M\) yield the same decoded
area. No global bound on \(f\) or recognition of the
integration-domain promise has been assumed.
\end{proof}
\end{annexitem}

\begin{annexitem}{annex:p04-fo-differentiation}{thm:p04-fo-differentiation}{A finite-prefix derivative obstruction}
\begin{proof}
The source functions \(f_n(x)=\sin(n\pi x)/n\) satisfy
\(\lVert f_n\rVert_\infty\leq1/n\), \(f_n'(0)=\pi\),
and \(f_n\in C^1[0,1]\).
For an entirely rational-polynomial alternative use
\(h_m(x)=x(1-x)^m\), \(m\geq1\).
Its derivative is
\((1-x)^{m-1}(1-(m+1)x)\); its maximum on \([0,1]\)
occurs at \(1/(m+1)\), with value at most \(1/(m+1)\).
Thus \(\lVert h_m\rVert_\infty\to0\), but \(h_m'(0)=1\).

Let \(p\) be the constant zero-polygon name of zero.
A continuous realizer for the scalar derivative at zero
would output a Cauchy name \(q\) of zero at \(p\), so
\(|q_3|\leq1/8\).
Continuity fixes that third scalar approximation on some
finite prefix of \(k\) zero-polygon entries.
Choose \(m\) with \(1/(m+1)<2^{-k}\).
The function \(h_m\) has a valid name beginning with that
prefix: those initial zero polygons meet their error bounds,
and subsequent polygons can be chosen by density.
It is in the differentiability promise domain.
The output's third approximation is still \(q_3\), but
it must now be within \(1/8\) of \(h_m'(0)=1\), impossible.
The same finite-prefix argument with sufficiently large
\(n\) and the source sine family gives its native obstruction.

If a continuous function-output derivative realizer existed,
compose it with the computable evaluation-at-zero realizer
already proved. This would give the forbidden scalar
realizer. Finite-use continuity consequently also rules out
computable derivative realizers.
Only uniform operator failure is proved by this argument;
no claim about one particular noncomputable derivative follows.
\end{proof}
\end{annexitem}

\begin{annexitem}{annex:p04-fo-derivative-promise}{thm:p04-fo-derivative-promise}{Secants under the declared derivative modulus}
\begin{proof}
For a consistent pair name of \(f,f'\), projecting the second
component directly yields a name of the derivative.
Now consider the different input consisting of \(f\)'s name
and the supplied valid derivative modulus \(\mu'\).
Fix \(n\), put \(\epsilon=2^{-n-2}\), choose
\(r\geq\mu'(n+2)\), and let \(h=2^{-r}\), \(N=2^r\).
Read \(P_k\), \(k=n+r+3\), whose uniform error is
\(\delta=2^{-n-r-3}\).
For \(j<N\) use
\[
 d_j=\frac{P_k((j+1)h)-P_k(jh)}h,
 \qquad d_N=d_{N-1}.
\]
These are effective rational values.

The corresponding exact secant of \(f\) is the average of
the continuous \(f'\) on that cell. Every point in the cell
is within \(h\) of its left endpoint, so the modulus bounds
its difference from \(f'(jh)\) by \(\epsilon\).
At the final node use the backward cell and the same bound.
Replacing the two function values by \(P_k\) changes a
secant by at most \(2\delta/h=\epsilon\).
Thus \(|d_j-f'(jh)|\leq2\epsilon\) at every node.

Interpolate the \(d_j\) linearly. At any point of a cell
the two true endpoint derivatives differ from the derivative
at that point by at most \(\epsilon\), again by the supplied
modulus. Convex interpolation therefore has error at most
\(3\epsilon<2^{-n}\).
Every finite output stage terminates using finitely much of
the two input streams, so it produces a derivative name.
The modulus belongs to \(f'\), not merely to \(f\);
none of its validity or availability was inferred from the
ordinary input name.
\end{proof}
\end{annexitem}

\begin{annexitem}{annex:p04-fo-general-zero}{thm:p04-fo-general-zero}{Plateau perturbations defeat even a relation realizer}
\begin{proof}
Let \(G_t\) be the source polygon with vertices
\((0,-1),(1/3,t),(2/3,t),(1,1)\).
It depends on \(t\) with
\(\lVert G_t-G_s\rVert_\infty\leq|t-s|\).
Let \(p\) repeat the rational polygon \(G_0\) at every position.
Suppose a continuous zero-relation realizer exists.
Its third scalar output approximation at \(p\) is some rational
\(q_3\). Continuity fixes that entry on a finite prefix of
\(k\) copies of \(G_0\).

Choose rational \(0<\epsilon<\min(1,2^{-k-1})\).
Both \(G_\epsilon\) and \(G_{-\epsilon}\) have valid names
beginning with that prefix, followed by their exact rational
polygon codes. Both have strict negative and positive
endpoint values and exactly one zero.
The former zero is strictly left of \(1/3\), the latter
strictly right of \(2/3\). Their distance exceeds \(1/3\).
The realizer must use the same \(q_3\) on both names,
within \(1/8\) of both zeros. That would put their distance
at most \(1/4\), a contradiction.
Thus even nonextensional relation realization fails on the
general domain, already with strict endpoint signs.
\end{proof}
\end{annexitem}

\begin{annexitem}{annex:p04-fo-nonextensional}{thm:p04-fo-nonextensional}{Native shrinking brackets and the selector counterfamily}
\begin{proof}
For a promised \(f\), continuity and rational density give
rational \(a_0<b_0\) with opposite strict signs.
Strict signs are semidecidable by evaluating the polygon
approximants and testing a value beyond its error bound.
Dovetail rational candidates to find the first certified pair.
Let \(\sigma\) be the sign at \(a_0\).

We verify the source's interior shrinking-pair observation.
Suppose \(\sigma f(a)>0>\sigma f(b)\).
Choose interior points with these signs near the endpoints.
The supremum \(c\) of the positive points between them
lies strictly before the negative endpoint and satisfies
\(f(c)=0\) by continuity. Positive points approach \(c\)
from the left. There are negative points arbitrarily close
on its right: after this supremum no positive point occurs,
and an interval containing neither sign would be a forbidden
zero interval. Strict-sign open neighbourhoods and rational
density therefore give \(a<u<v<b\), as close to \(c\)
as needed, with \(\sigma f(u)>0>\sigma f(v)\).
In particular \(v-u\leq(b-a)/2\).

Enumerate such rational pairs and dovetail their strict-sign
tests. The observation ensures termination at every stage.
The intervals are nested, have width at most \(2^{-j}\)
at stage \(j\), and maintain their endpoint signs.
Their endpoints converge to a common \(z\in[0,1]\).
Continuity gives both \(\sigma f(z)\geq0\) and
\(\sigma f(z)\leq0\), hence \(f(z)=0\).
Outputting the left endpoints gives host error at most
\(2^{-j}\). It also satisfies the native strict pairwise
scalar clause: later left endpoints lie strictly between
the current endpoints, so their difference is less than
the current width, at most \(2^{-j}\).
All searches are name-based effective procedures.

An extensional continuous selector is nevertheless impossible.
Use the source family \(H_t\) with vertices
\((0,-1),(1/3,t),(2/3,t-1),(1,1)\).
Every member has endpoint sign change and no zero interval.
Its middle slope is fixed; any constant outer segment has
nonzero value. From a Cauchy name \(q_n\) of \(t\), the
polygons \(H_{q_n}\) form a computable host name of \(H_t\),
since their uniform error is at most \(|q_n-t|\).

If an extensional continuously realized selector \(Z\)
existed, \(r(t)=Z(H_t)\) would be a total continuously
realized scalar map. The scalar quotient/continuity theorem
from the foundation makes \(r\) topologically continuous.
At \(t=2\) the unique zero is \(1/9<1/3\).
On the connected interval \((1/3,3)\), \(H_t(1/3)=t\ne0\);
continuity prevents \(r\) from crossing \(1/3\).
Thus \(r(t)<1/3\) throughout that interval.
At \(t=-1\) the unique zero is \(8/9>2/3\).
On \((-2,2/3)\), \(H_t(2/3)=t-1\ne0\), so similarly
\(r(t)>2/3\) throughout. The two parameter intervals
overlap, a contradiction. This proves the selector clause
without treating relation outputs as extensional.
\end{proof}
\end{annexitem}

\begin{annexitem}{annex:p04-fo-unique-zero}{thm:p04-fo-unique-zero}{Compact exclusion, endpoints and supplied local isolation}
\begin{proof}
For rational \(u<v\), the complement
\(K=[0,1]\setminus(u,v)\) is a finite union of rational
compact intervals, or empty.
Its polygon minimum absolute value is effective as described
in the naming proof. If \(\min_K|P_k|>2^{-k}\), then
\(|f(x)|\geq|P_k(x)|-2^{-k}>0\) for all \(x\in K\).
Every zero consequently belongs to \((u,v)\).

Suppose the sole zero is \(z\). For each precision choose
rational \(u<z<v\) with \(v-u<2^{1-n}\), allowing
endpoints outside \([0,1]\).
If \(K\) is nonempty, continuity and compactness give
\(\eta=\min_K|f|>0\). For sufficiently large \(k\),
\(\min_K|P_k|\geq\eta-2^{-k}>2^{-k}\).
If \(K\) is empty the test is vacuous.
Enumerating rational \(u,v\) and integers \(k\) and testing
these conditions therefore terminates at every output position.
The midpoint has distance less than \(2^{-n}\) from \(z\).
This is a productive scalar realizer on all unique-zero names.
The source's stronger margin \(2\cdot2^{-k}\) is also
eventually met, by choosing \(3\cdot2^{-k}<\eta\);
it changes the test, not the conclusion.

For a supplied rational \(0\leq a<b\leq1\), restrict and
rescale to \(g(t)=f(a+(b-a)t)\).
A polygon restricts/rescales by rational breakpoints and
has exactly the same error bound, so this operation on
names is effective. If the supplied interval has exactly
one zero, apply the preceding algorithm and return
\(a+(b-a)z\) by scalar arithmetic.
If \(a=b\) and its endpoint is promised a zero, return
the rational endpoint without testing that equality.
Zeros elsewhere do not affect this supplied local promise.
An isolated interior zero has some such rational isolating
interval. For one fixed computable function it can be
hard-coded, giving individual computability; this does not
give an algorithm for finding the interval uniformly.
Endpoint isolated zeros are already the rational zero or one.
\end{proof}
\end{annexitem}

\begin{annexitem}{annex:p04-fo-inverse}{thm:p04-fo-inverse}{Global inverse values and names}
\begin{proof}
Global evaluation is computable: bound a scalar input using
its zeroth approximation, restrict the global boxes to a
large rational compact interval, rescale to \([0,1]\),
and use the naming and evaluation algorithms.
For strictly increasing onto \(f\), each \(y\) has a unique
preimage \(z\). At precision \(n\), search rational \(a<b\)
with \(b-a<2^{-n}\), \(f(a)<y<f(b)\).
The strict comparisons are semidecidable from scalar names.
Continuity and strict monotonicity supply such witnesses
arbitrarily near \(z\), so the dovetailed search terminates.
Their midpoint is a valid scalar approximation to \(z\).

For a name of the entire inverse, enumerate ordered rational
domain intervals \([c,d]\) and rational \(a<b\).
Strict monotonicity implies exactly
\[
 f^{-1}([c,d])\subseteq(a,b)
 \quad\Longleftrightarrow\quad f(a)<c\ \text{and}\ d<f(b).
\]
Here the left side denotes the image of the interval under
the inverse function. Both strict tests are semidecidable,
so dovetailing them enumerates all and only its proper
box properties. Empty-domain source boxes are restored
mechanically. This is a global box-name of the inverse.
Neither strictness nor surjectivity has been tested; both
are exactly the mathematical input-domain promise.
\end{proof}
\end{annexitem}

\begin{annexitem}{annex:p04-fo-approximate-zero}{thm:p04-fo-approximate-zero}{The certified residual search}
\begin{proof}
Enumerate rational \(s\in[0,1]\) and precision indices \(k\),
and test \(|P_k(s)|+2^{-k}<2^{-n}\).
Any successful pair guarantees
 \(|f(s)|\leq|P_k(s)|+2^{-k}<2^{-n}\).
If a rational \(s\) satisfies the required strict inequality,
choose \(k\) with \(2^{1-k}<2^{-n}-|f(s)|\);
then the certificate passes.
If only a real point with this strict residual is known to
exist, continuity and rational density give a rational one
in \([0,1]\); a root at zero or one itself is rational.
Thus the dovetailed search terminates on the stated domain.
At a coarse precision a small \(P_k(s)\) alone would give
no such guarantee: the missing error term cannot be dropped.
\end{proof}
\end{annexitem}

\begin{annexitem}{annex:p04-fo-no-computable-zero}{thm:p04-fo-no-computable-zero}{A complete small-cover and tent-function construction}
\begin{proof}
Use the fixed effective numerical program listing
\(\varphi_e\) from \cref{thm:enumeration}.
For each \(e\), dovetail the computation
\(\varphi_e(e+5)\). If it converges with a natural output,
decode that output by \(\nu_{\mathbb Q}\) as \(q_e\) and
enumerate the open interval of centre \(q_e\) and radius
\(2^{-e-4}\). Its length is \(2^{-e-3}\).
Every Cauchy-computable real \(x\) is covered: its rational
code sequence is a total computable numerical function,
has some index \(e\), and its \((e+5)\)-th approximation
has error at most \(2^{-e-5}\), strictly less than that
interval radius. Invalid or partial programs do not impair
the construction; each index contributes at most one interval.

The sum of all possible lengths is \(1/4\).
The resulting open set \(U\) cannot cover \([0,1]\):
otherwise compactness gives a finite interval subcover,
whose total interval length must be at least one, contrary
to the budget. The finite-length fact follows by sorting
all endpoints and summing the covered elementary segments;
each segment's length is counted at least once by the
cover's interval lengths.
Thus \([0,1]\setminus U\) is nonempty, and none of its
points is Cauchy-computable.

Turn the effective interval enumeration into a total
stage sequence, allowing empty padding intervals.
For a proper interval \(I_j=(a_j,b_j)\), define
\[
 T_j(x)=\max(0,\min(1,x-a_j,b_j-x));
\]
padding has \(T_j=0\).
Each restriction to \([0,1]\) is an effectively coded
rational polygon, has \(0\leq T_j\leq1\), and is positive
exactly on \(I_j\cap[0,1]\).
Its breakpoints are rational: the interval endpoints,
their midpoint, and, when a unit-height plateau occurs,
\(a_j+1,b_j-1\), clipped to the domain.
Let
\[
 F(x)=\sum_{j=0}^\infty 2^{-j-1}T_j(x).
\]
The partial sum of the first \(N\) terms is a rational
polygon obtained on the common finite breakpoint grid.
The remaining uniform error is at most
\(\sum_{j\geq N}2^{-j-1}=2^{-N}\).
These effective partial sums are a host name of a
continuous \(F\). Nonnegative terms show \(F(x)>0\)
exactly when \(x\in U\), so its zero set is the nonempty
complement, with no computable point.
Finally \(-F\) is computable, has maximum value zero,
and its maximum points are exactly those zeros.
This supplies the second clause and keeps maximum value
distinct from maximum-point selection.
\end{proof}
\end{annexitem}
